\subsection{Factoring with Rational Roots}
\begin{Example}
Use the rational roots theorem to find at least one factor of the polynomial
\[
4x^4-8x^3-7x^2+11x+6
\]
First, we find all possible roots.
\begin{itemize}
\item $a_n=\makeblank[.5in]$, so the possible values of $q$ are \makeblank
\item $a_0=\makeblank[.5in]$, so the possible values of $p$ are \makeblank
\end{itemize}
\begin{lpic}{}{4}
\foreach \y/\val in {0/1,1/2,2/4}
  \regtextleft{1}{-2-\y}{$q=\pm\val$};
\foreach \x/\val in {0/1,1/2,2/3,3/6}
  \regtext{3+\value{cols}/4-3/8+\x*\value{cols}/2-3*\x/4}{-1}{$p=\pm\val$};
\end{lpic}
The possible rational roots are \makeblank[3in]. We will try numbers from this list until we find one that is a factor. The easiest way is to use synthetic division.
\begin{cpic}{}{7}
\foreach \y in {0,-4} {
	\synthdiv{-6}{\y}{5};
	\foreach \x/\lbl in {0/4,1/-8,2/-7,3/11,4/6}
		\regtext{-3.5+2*\x}{\y-1}{$\lbl$};
}
\end{cpic}
We got lucky and found a factor in only two tries.
\[
4x^4-8x^3-7x^2+11x+6=\Bigl(\makeblank[1in]\Bigr)\Bigl(\makeblank\Bigr)
\]
Notice that the degree of each factor is \makeblank[1in] than the degree of the original. The new polynomials are ``\makeblank[1in].''
\end{Example}
Note: We could have ended up with different factors. See what happens if you try $r=-1$ instead of $r=2$.